%% kapa: exam=1 round=21 year=2010 period=1 form=A subject=civil_law subject_ko=민법 seq=27 number=27
%% source: Q-Net 감정평가사 기출문제 (cst00309) artlSeq=5024181 file=extracted/1008/5024181/2010년 제21회 감정평가사 1차시험 기출문제/1교시.hwp
%% conversion: kordoc 4.24.51 -> tools/gen1.py (auto); review=proofread-sonnet; flags=none
\begin{question}{27}
불공정한 법률행위에 관한 설명으로 옳지 않은 것은? (다툼이 있으면 판례에 의함)
\begin{choices}
  \item 무효인 불공정한 법률행위는 추인에 의하여 유효로 될 수 있다.
  \item 불공정한 법률행위가 되기 위해 궁박, 경솔, 무경험 세 가지 모두가 요구되는 것은 아니다.
  \item 불공정한 법률행위의 요건에 대한 증명책임은 그 무효를 주장하는 자에게 있다.
  \item 궁박은 경제적 궁박에 한정되지 않으며, 정신적 또는 심리적 원인에 기한 궁박도 포함된다.
  \item 대리인에 의한 법률행위에서 경솔․무경험은 대리인을 기준으로 판단하지만, 궁박 여부는 본인을 기준으로 판단한다.
\end{choices}
\end{question}
\answer{1}
